In this Appendix, I present a synthetic example where the relationship between covariates $Z_g$ and the group-specific propensities $\beta_g$ is non-linear. 

The purpose of this example is to look at the differential performance of \texttt{seine}, which is explicitly designed to capture non-linear relationships, and a linear estimator, in estimating non-linear relationships and, finally, the associated group-specific propensities.

\subsection{Data generating process}

The setup is similar to the one presented in Section \ref{sec:A synthetic example}, with the difference that now we do not have a regional distinction, but rather a continuous covariate $Z_g$.

Let us suppose that $Z_g$ represents the share of migrants within a precinct, which implies $Z_g$ continuous and $Z_g\in [0,1]$.

The data generating process is the following.

Precincts differ for their percentage of high-income and low-income voters: the percentages of high-income voters $\overline{X}_g^H$ are extracted from a truncated normal distribution (truncated at 0 and 1) with mean 0.5 and standard deviation 0.2:

\[
\overline{X}_g^H \sim \text{TruncNormal}(a=0, b=1, \mu=0.5, \sigma=0.2).
\]

We assume that the share of immigrants $Z_g$ has a loose negative association with the share of high-income voters $\overline{X}_g^H$ (migrants are more likely to live in low-income areas), creating possible identification issues.

We can model this association by assuming that the share of immigrants is extracted from a truncated normal distribution with mean $1-\overline{X}_g^H$ and standard deviation 0.2:

\[
Z_g \sim \text{TruncNormal}(a=0, b=1, \mu=1-\overline{X}_g^H, \sigma=0.2).
\]

Then, we might assume that

\begin{align*}
\beta_{g}^{H} &\sim \text{TruncNormal}(a=0, b=1, \mu=\eta_{g}^{H}, \sigma=0.05), \\
\beta_{g}^{L} &\sim \text{TruncNormal}(a=0, b=1, \mu=\eta_{g}^{L}, \sigma=0.05).
\end{align*}

The crucial difference with previous setup arises when we set means to be non-linear functions of $Z_g$:

\begin{align*}
\eta_{g}^{H} = 10^{Z_g-1}, \quad & \eta_{g}^{L} = e^{Z_g-1}.
\end{align*}

The interpretation of this condition is that both high-income and low-income voters tend to favor the reform more as the share of immigrants in their precinct increases, but the effect is stronger for low-income voters and the relationship is non-linear.

Finally, the observed precinct-level outcome ($\overline{Y}_{g}$, i.e., the share of voters in favour of the reform in precinct $g$) is generated as the weighted average of these two group-specific propensities:
\[
\overline{Y}_{g}=\overline{X}_{g}^{H}\beta_{g}^{H}+\overline{X}_{g}^{L}\beta_{g}^{L}.
\]

\subsection{Comparison of results}

The true values of the group-specific propensities are $\beta_{g}^{H} = 0.324$ and $\beta_{g}^{L} = 0.660$.

% latex table generated in R 4.5.1 by xtable 1.8-4 package
% Thu Sep 10 13:01:24 2026
\begin{table}[ht]
\centering
\begin{tabular}{rr}
  \hline
high & low \\ 
  \hline
0.324 & 0.660 \\ 
   \hline
\end{tabular}
\caption{True coefficients up to approximation} 
\label{tab:true-syn-exp-beta}
\end{table}


\subsubsection{OLS}

% latex table generated in R 4.5.1 by xtable 1.8-4 package
% Thu Sep 10 13:01:24 2026
\begin{table}[ht]
\centering
\begin{tabular}{rrrrr}
  \hline
 & Estimate & Std. Error & t value & Pr(>|t|) \\ 
  \hline
(Intercept) & 0.917 & 0.023 & 40.423 & 0.000 \\ 
  high\_income & -0.845 & 0.042 & -20.021 & 0.000 \\ 
   \hline
\end{tabular}
\caption{Coefficients of naive OLS regression} 
\label{tab:naive-regression-exp}
\end{table}


OLS estimates are $\hat{\beta}_{g}^{H} = 0.917$ and $\hat{\beta}_{g}^{L} = 0.917-0.845 = 0.072$.

As we could expect, these estimates are far from the true values as they do not take into account the effect of the covariate $Z_g$ on the group-specific propensities $\beta_g$.

\subsubsection{\texttt{seine}}

% latex table generated in R 4.5.1 by xtable 1.8-4 package
% Wed Sep  9 17:52:51 2026
\begin{table}[ht]
\centering
\begin{tabular}{llrrrr}
  \hline
predictor & outcome & estimate & std.error & conf.low & conf.high \\ 
  \hline
high\_income & outcome & 0.321 & 0.014 & 0.294 & 0.348 \\ 
  low\_income & outcome & 0.664 & 0.023 & 0.618 & 0.709 \\ 
   \hline
\end{tabular}
\caption{Semiparametric ecological inference estimates} 
\label{tab:ei-estimates-exp}
\end{table}


As in the example presented in Section \ref{sec:A synthetic example}, \texttt{seine} is able to provide good estimates of the group-specific propensities.

\subsubsection{Interaction between $\overline{X}_g$ and $Z_g$}

% latex table generated in R 4.5.1 by xtable 1.8-4 package
% Thu Sep 10 13:01:24 2026
\begin{table}[ht]
\centering
\begin{tabular}{rrrrr}
  \hline
 & Estimate & Std. Error & t value & Pr(>|t|) \\ 
  \hline
(Intercept) & 0.176 & 0.022 & 7.866 & 0.000 \\ 
  high\_income & -0.082 & 0.034 & -2.403 & 0.017 \\ 
  zeta & 0.869 & 0.033 & 26.049 & 0.000 \\ 
  high\_income:zeta & -0.357 & 0.060 & -5.961 & 0.000 \\ 
   \hline
\end{tabular}
\caption{Coefficients of regression with interactions with covariates} 
\label{tab:interaction-regression-exp}
\end{table}


In this case, the functional form of the relationship between $Z_g$ and $\beta_g$ is estimated to be linear, which is not the case in the data generating process. In particular, it is estimated that $E[\beta_g^L | Z_g] = \gamma_0^L + \gamma_1^L*Z_g$, while $E[\beta_g^H | Z_g] = \gamma_0^H + \gamma_1^H*Z_g$.

The estimates are $\hat{\beta}_{g}^{L} = 0.176 + (0.869)*Z_g$ and $\hat{\beta}_{g}^{H} = (0.176-0.082) + (0.869-0.357)*Z_g$.

Then, we can adopt the same strategy presented in Section 4 to produce global estimates.

If we plug in the observed values of $Z_g$ and $\overline{X_g}$ in the above equations, we obtain estimates of $\beta_g$ for each precinct $g$.
The weighted average of these estimates across precincts is $\hat{\beta}^{H} = 0.321$ and $\hat{\beta}^{L} = 0.664$, which are in fact quite close to the true values of $\beta^{H}$ and $\beta^{L}$ and correspond to estimates performed with \texttt{seine}.

% latex table generated in R 4.5.1 by xtable 1.8-4 package
% Wed Sep  9 17:52:51 2026
\begin{table}[ht]
\centering
\begin{tabular}{rr}
  \hline
beta\_high\_hat & beta\_low\_hat \\ 
  \hline
0.321 & 0.664 \\ 
   \hline
\end{tabular}
\caption{Estimates with covariates} 
\label{tab:int-syn-exp-beta}
\end{table}


The conclusion that we can draw from this example is that, even if the relationship between $Z_g$ and $\beta_g$ is non-linear, a linear approximation of this relationship can still provide good estimates of the group-specific propensities $\beta_g$, as far as the relevant covariates are used.